% =========================================================
% doc2-polar-semiring.tex — Section 2: Primitive Polarity on Semirings
% =========================================================

\section{Primitive Polarity on Semirings}\label{sec:ps}

The objective of this section is to establish primitive polarity as an
independent algebraic primitive that can be adjoined to an idempotent
commutative semiring without appeal to residuals or dual operations. We
first fix the underlying semiring, then introduce the polarity through two
further axioms, isolate the role of each, exhibit canonical models, and
verify that the resulting eleven-axiom system is independent.

\begin{definition}[Primitive semiring]\label{def:semiring}
A \emph{primitive semiring} is a structure $(X,\vee,\otimes,\ev,\eo)$, with
$\vee,\otimes$ binary operations and $\ev,\eo \in X$, satisfying for all
$a,b,c \in X$:
\begin{align}
  a \vee a &= a, \tag{J0}\label{ax:J0}\\
  (a \vee b) \vee c &= a \vee (b \vee c), \tag{J1}\label{ax:J1}\\
  a \vee b &= b \vee a, \tag{J2}\label{ax:J2}\\
  a \vee \ev &= a, \tag{J3}\label{ax:J3}\\
  (a \otimes b) \otimes c &= a \otimes (b \otimes c), \tag{T1}\label{ax:T1}\\
  a \otimes b &= b \otimes a, \tag{T2}\label{ax:T2}\\
  a \otimes \eo &= a, \tag{T3}\label{ax:T3}\\
  a \otimes (b \vee c) &= (a \otimes b) \vee (a \otimes c),
    \tag{TJ1}\label{ax:TJ1}\\
  a \otimes \ev &= \ev. \tag{TJ2}\label{ax:TJ2}
\end{align}
\end{definition}

Axioms \eqref{ax:J0} to \eqref{ax:J3} make $(X,\vee,\ev)$ an idempotent
commutative monoid, axioms \eqref{ax:T1} to \eqref{ax:T3} make
$(X,\otimes,\eo)$ a commutative monoid, and \eqref{ax:TJ1} to \eqref{ax:TJ2}
link the two, so that $(X,\vee,\otimes,\ev,\eo)$ is an idempotent commutative
semiring \citep{Golan1999}. These nine axioms describe the algebraic layer
alone; no duality, reflection, or residual is present.

\begin{definition}[Polar Semiring]\label{def:ps}
A \emph{Polar Semiring} is a structure $(X,\vee,\otimes,{}^{*},\ev,\eo)$
whose reduct $(X,\vee,\otimes,\ev,\eo)$ is a primitive semiring and whose
unary operation ${}^{*}\colon X \to X$ satisfies, for all $a,b \in X$:
\begin{align}
  \dstar{(\dstar{a})} &= a, \tag{P}\label{ax:P}\\
  a \vee b = b \;&\Longleftrightarrow\; \dstar{b} \vee \dstar{a} = \dstar{a}.
    \tag{PJ}\label{ax:PJ}
\end{align}
The operation ${}^{*}$ is the \emph{polarity} of the structure. We
abbreviate ``Polar Semiring'' by PS.
\end{definition}

% Axioms \eqref{ax:P} and \eqref{ax:PJ} are of a different nature from the
% nine semiring axioms: they add no algebraic operation but introduce a
% primitive polarity and its single point of contact with the semiring. To
% make that contact explicit, define the \emph{join order} by
% \begin{equation}\label{eq:orders}
%   a \leq b \;:\Longleftrightarrow\; a \vee b = b,
% \end{equation}
% which is a partial order by \eqref{ax:J0} to \eqref{ax:J2}, with least
% element $\ev$ by \eqref{ax:J3} and $a \vee b$ the join of $a$ and $b$. Under
% \eqref{eq:orders}, axiom \eqref{ax:PJ} states that $a \leq b$ if and only if
% $\dstar{b} \leq \dstar{a}$; that is, ${}^{*}$ is an antitone involution on
% $(X,\leq)$, the classical notion of a polarity on an ordered set
% \citep{Birkhoff1940}. No dual operation belongs to the signature, and none
% is constructed here; the structures generated from ${}^{*}$ by conjugation
% are the subject of \Cref{sec:arch}. The next two propositions record the
% distinct roles of \eqref{ax:P} and \eqref{ax:PJ}.



\begin{proposition}[Role of the polarity axiom]\label{prop:role-P}
Under \eqref{ax:P}, the polarity ${}^{*}$ is a bijection with
${}^{*} = {}^{*}{}^{-1}$.
\end{proposition}

\begin{proof}
By \eqref{ax:P} each $a$ equals $\dstar{(\dstar{a})}$, so ${}^{*}$ is
surjective; and $\dstar{a} = \dstar{b}$ gives
$a = \dstar{(\dstar{a})} = \dstar{(\dstar{b})} = b$, so ${}^{*}$ is
injective. The same equation shows ${}^{*}$ equals its own inverse.
\end{proof}

\begin{proposition}[Nontriviality of polarity]\label{prop:role-PJ}
Let $(X,\vee,{}^{*})$ satisfy \eqref{ax:J0} to \eqref{ax:J2} and
\eqref{ax:PJ}. If ${}^{*} = \mathrm{id}_X$, then $X$ is a singleton.
Consequently, on every nontrivial Polar Semiring, ${}^{*} \neq \mathrm{id}_X$.
\end{proposition}

\begin{proof}
Assume ${}^{*} = \mathrm{id}_X$, let $a,b \in X$ be arbitrary, and set
$c = a \vee b$. By \eqref{ax:J1} and \eqref{ax:J0},
$$
  a \vee c = a \vee (a \vee b) = (a \vee a) \vee b = a \vee b = c.
$$
Applying \eqref{ax:PJ} to the pair $(a,c)$, the equality $a \vee c = c$
yields $\dstar{c} \vee \dstar{a} = \dstar{a}$, which under
${}^{*} = \mathrm{id}_X$ reads $c \vee a = a$; but \eqref{ax:J2} gives
$c \vee a = a \vee c = c$, so $a = c$. Symmetrically, by \eqref{ax:J1},
\eqref{ax:J0}, and \eqref{ax:J2},
$$
  b \vee c = b \vee (a \vee b) = (b \vee b) \vee a = b \vee a = a \vee b = c,
$$
and \eqref{ax:PJ} applied to $(b,c)$ gives $\dstar{c} \vee \dstar{b} =
\dstar{b}$, that is, $c \vee b = b$; again \eqref{ax:J2} gives $c \vee b =
b \vee c = c$, so $b = c$. Hence $a = c = b$. As $a,b$ were arbitrary, any
two elements of $X$ coincide, so $X$ is a singleton.
\end{proof}

The proof uses only the join idempotence, associativity, and commutativity
\eqref{ax:J0} to \eqref{ax:J2} together with \eqref{ax:PJ}; neither the join
unit \eqref{ax:J3}, nor any tensor axiom, nor the order or the derived meet
is invoked. Thus \eqref{ax:PJ} does more than relate the primitive semiring
to the involution: on every nontrivial model it forbids the trivial
involution, so that the polarity of a nontrivial Polar Semiring is a genuine
reflection rather than the identity. Axiom \eqref{ax:P} alone would permit
${}^{*} = \mathrm{id}_X$; it is \eqref{ax:PJ} that rules this out. The deeper
structural role of \eqref{ax:PJ}, generating the compatibility between a
Polar Semiring and its dual, is established in \Cref{sec:arch}.

We now exhibit canonical models, confirming that the class is nonempty and
that primitive polarity coexists with the semiring operations. In each case
the join order coincides with the natural order of a chain, so
\eqref{ax:J0} to \eqref{ax:J3} are immediate; we verify the remaining
axioms directly.

\begin{example}[Involutive semiring model]\label{ex:involutive-semiring}
An \emph{involutive semiring} consists of a structure
$(X,\vee,\otimes,{}^{*},e_{\vee},e_{\otimes})$,
where $(X,\vee,e_{\vee})$ and $(X,\otimes,e_{\otimes})$ are commutative monoids,
${}^{*}:X\to X$ is an involution satisfying
$(a^{*})^{*}=a$,
and distributes over both operations,
\[
(a\vee b)^{*}=a^{*}\vee b^{*},\qquad
(a\otimes b)^{*}=a^{*}\otimes b^{*}.
\]
Moreover, $\otimes$ distributes over $\vee$.

Compared with Polar Semirings, this structure does not require
$\vee$ to be idempotent and does not impose any compatibility between
${}^{*}$ and the order induced by $\vee$. In particular, the identity map
${}^{*}=\mathrm{id}_{X}$ always satisfies the involution and distribution
axioms. Hence involutive semirings illustrate that the semiring axioms
together with involution alone do not exclude the trivial involution,
highlighting the essential role of the polarity axiom \eqref{ax:PJ}.
\end{example}

\begin{example}[Boolean model]\label{ex:boolean}
Let $X = \{0,1\}$ with $\vee = \max$, $\otimes = \min$, $\dstar{x} = 1-x$,
$\ev = 0$, $\eo = 1$. The monoid axioms \eqref{ax:T1} to \eqref{ax:T3} hold
for $\min$ with unit $1$; \eqref{ax:TJ1} is the distributivity of $\min$
over $\max$; \eqref{ax:TJ2} reads $\min(a,0)=0$; \eqref{ax:P} holds since
$1-(1-x)=x$; and \eqref{ax:PJ} holds because $x \mapsto 1-x$ is antitone.
This is the smallest nontrivial Polar Semiring.
\end{example}

\begin{example}[Tropical model]\label{ex:tropical}
Let $X = \mathbb{R} \cup \{-\infty,+\infty\}$ with $\vee = \max$,
$a \otimes b = a+b$, $\dstar{x} = -x$, $\ev = -\infty$, $\eo = 0$, under the
convention $(-\infty)+x = -\infty$ for all $x$, including $x = +\infty$. Then
$(X,\otimes,0)$ is a commutative monoid, so \eqref{ax:T1} to \eqref{ax:T3}
hold; addition is monotone, so $a + \max(b,c) = \max(a+b,a+c)$ gives
\eqref{ax:TJ1}; the stated convention gives $a \otimes (-\infty) = -\infty$,
which is \eqref{ax:TJ2}; \eqref{ax:P} holds since $-(-x)=x$; and negation is
antitone, giving \eqref{ax:PJ}.
\end{example}

\begin{example}[Quadratic model on the unit interval]\label{ex:quadratic}
Let $X = [0,1]$ with $\vee = \max$, $a \otimes b = ab$, and
\begin{equation}\label{eq:quadratic-star}
  \dstar{x} := \bigl(1 - \sqrt{x}\bigr)^{2},
\end{equation}
with $\ev = 0$ and $\eo = 1$. Ordinary multiplication makes
$(X,\otimes,1)$ a commutative monoid and is monotone, giving \eqref{ax:T1}
to \eqref{ax:T3} and \eqref{ax:TJ1}; $a \otimes 0 = 0$ gives \eqref{ax:TJ2}.
The map \eqref{eq:quadratic-star} is strictly decreasing with $\dstar{0}=1$,
$\dstar{1}=0$, whence \eqref{ax:PJ}, and it is an involution: since
$1 - \sqrt{x} \geq 0$ on $[0,1]$,
$$
  \dstar{(\dstar{x})}
  = \Bigl(1 - \sqrt{(1-\sqrt{x})^{2}}\Bigr)^{2}
  = \bigl(1 - (1 - \sqrt{x})\bigr)^{2}
  = x,
$$
which is \eqref{ax:P}. This model plays a decisive role in
\Cref{subsec:impossibility}.
\end{example}

No dual semiring is introduced or analyzed at this stage; the three models
are given purely in the primitive signature $(X,\vee,\otimes,{}^{*},\ev,\eo)$.
We close by verifying that no axiom is redundant.

\begin{theorem}[Independence]\label{thm:independence}
The eleven-axiom system of \Cref{def:semiring} and \Cref{def:ps} is
independent: for each axiom there exists a finite structure
$(X,\vee,\otimes,{}^{*},\ev,\eo)$ satisfying the remaining ten axioms but not
the selected one.
\end{theorem}

\begin{proof}
For each axiom a finite countermodel was obtained by exhaustive search with
the model finder Mace4, and the nonderivability claims were checked
independently with the theorem prover Prover9 \citep{McCune2005}. The
countermodels have at most \noteAI{fill in the maximal cardinality} elements
and are listed in \noteAI{decide: appendix with the model tables, or
supplementary material}.
\noteAI{Per final rules \S6.2, the original Mace4/Prover9 runs used the
obsolete axiom names (M0 to M3, S1 to S3, SM1 to SM2, PM) and the obsolete
choice of primitives ($\wedge$ primitive). The runs must be repeated in the
present signature (primitives $\vee,\otimes,{}^{*},\ev,\eo$; axioms J0 to J3,
T1 to T3, TJ1, TJ2, P, PJ) before this theorem is submitted.}
\end{proof}

Each axiom therefore contributes an essential and independent ingredient to
the definition of a Polar Semiring. This section has shown that primitive
polarity may be adjoined to a semiring consistently and independently, yet
so far it is only an involution constrained by a single compatibility
axiom. Whether these minimal assumptions already determine a dual
architecture is the question of the next section, which answers it in the
affirmative: primitive polarity canonically generates a dual semiring
together with a compatible order.
